\documentclass[a4paper, 12pt, titlepage]{article}
% ==========================================
% 1. Pacotes básicos
% ==========================================
% Margens ABNT: Superior 3cm, Esquerda 3cm, Inferior 2cm, Direita 2cm
\usepackage[top=3cm, left=3cm, bottom=2cm, right=2cm]{geometry}

% Espaçamento entre linhas de 1,5
\usepackage{setspace}
\onehalfspacing

% Recuo do primeiro parágrafo (1,25 cm) e controle de espaçamento entre parágrafos
\usepackage{indentfirst}
\setlength{\parindent}{1.25cm}
\setlength{\parskip}{0pt}

% Idioma e fontes
\usepackage[brazil]{babel}
\usepackage[utf8]{inputenc}
\usepackage[T1]{fontenc}
\usepackage{times}

%  Matemática
\usepackage{amsmath}
\usepackage{amssymb}
\usepackage{amsthm}

% Imagens
\usepackage{graphicx}

% Sumário
\usepackage{hyperref}
\hypersetup{
    colorlinks=true,
    linkcolor=black,
    citecolor=black,
    urlcolor=black,
    filecolor=black,
    bookmarksnumbered=true,
    pdfstartview={FitH}
}

% ==========================================
% 2. AMBIENTE PARA CÓDIGO R 
% ==========================================
\usepackage{listings}
\usepackage{xcolor}

% Paleta de cores para realce de sintaxe do R
\definecolor{r_bg}{rgb}{0.97, 0.97, 0.98}        % Fundo suave
\definecolor{r_comment}{rgb}{0.40, 0.45, 0.55}   % Comentários
\definecolor{r_keyword}{rgb}{0.10, 0.30, 0.70}   % Palavras-chave
\definecolor{r_string}{rgb}{0.15, 0.55, 0.15}    % Strings
\definecolor{r_numbers}{rgb}{0.5, 0.5, 0.5}      % Numeração de linhas
\definecolor{r_rule}{rgb}{0.85, 0.85, 0.88}      % Borda do quadro

% Definição do ambiente \begin{rcode} ... \end{rcode}
\lstnewenvironment{rcode}[1][]{%
  \lstset{
    language=R,
    backgroundcolor=\color{r_bg},
    basicstyle=\ttfamily\small,
    keywordstyle=\color{r_keyword}\bfseries,
    commentstyle=\color{r_comment}\itshape,
    stringstyle=\color{r_string},
    numberstyle=\tiny\color{r_numbers},
    numbers=left,
    numbersep=10pt,
    stepnumber=1,
    showstringspaces=false,
    tabsize=2,
    breaklines=true,
    breakatwhitespace=true,
    frame=single,
    rulecolor=\color{r_rule},
    frameround=tttt,
    captionpos=b,
    xleftmargin=15pt,
    framexleftmargin=15pt,
    extendedchars=true,
    literate=
      {á}{{\'a}}1 {é}{{\'e}}1 {í}{{\'i}}1 {ó}{{\'o}}1 {ú}{{\'u}}1
      {Á}{{\'A}}1 {É}{{\'E}}1 {Í}{{\'I}}1 {Ó}{{\'O}}1 {Ú}{{\'U}}1
      {à}{{\`a}}1 {à}{{\`A}}1
      {ã}{{\~a}}1 {õ}{{\~o}}1 {Ã}{{\~A}}1 {Õ}{{\~O}}1
      {â}{{\^a}}1 {ê}{{\^e}}1 {ô}{{\^o}}1 {Â}{{\^A}}1 {Ê}{{\^E}}1 {Ô}{{\^O}}1
      {ç}{{\c{c}}}1 {Ç}{{\c{C}}}1,
    #1
  }
}{}

\title{HOMEWORK 1}
\author{Arthur Augusto Vieira Pinho Gimenez (579859)\\
        Caio Felipe dos Santos Oliveira  (516882) \\
        Roberto Walace Santos da Silva (581562)}
\date{\today}

\begin{document}

\maketitle
\tableofcontents
\newpage

\section{Questão 1}

A primeira questão solicita que uma amostra de 300 observações do conjunto de dados em anexo, armazenado em formato .csv, seja carregada pelo R e armazenada no objeto \texttt{data\_group}. No entanto, é notável que o arquivo \textit{HW1\_bike\_sharing} apresenta 731 observações, portanto precisamos extrair apenas 300 do total, seguindo o critério requisitado pela questão.

Considere $M$ como o maior número de matrícula entre os integrantes do grupo, ou seja, $M=max \, \{579859, 516882, 581562\}=581562$. Defina $r$ como a posição inicial das 300 observações, cuja fórmula é:
\begin{gather*}
    r= 1 +(M \bmod 100) \\
    \intertext{Substituindo o valor de $M$ na fórmula acima,} 
    r= 1 +(581562 \bmod 100) \\
    r= 63
\end{gather*}
Defina $f$ como a posição final dos dados extraídos, segue, por aritmética simples:
\begin{gather*}
    300= f-63 +1\\
    f=362
\end{gather*}

Com base nessas duas posições, é possível determinar a data da primeira e última observação no arquivo .csv, conforme solicitado na questão.

% Primeira tabela
\begin{table}[h!]
    \centering
    \caption{Linhas da posição 63 e 362}
    \label{tab:mascara_1u}
    \begin{tabular}{c|c|c|c|c|c|c|c}
        &instant &dteday &season &weathersit &temp &casual &registered \\ \hline
        63 &63 &2011-03-04 &1 &2 &10.7 &214 &1730 \\ \hline
        362 &362 &2011-12-28 &1 &1 &12.3 &255 &2047
    \end{tabular}
\end{table}

Portanto, ao analisar as linhas acima, é conclusivo que as datas referentes à observação 63 e à 362 são respectivamente \textit{2011-03-04} e \textit{2011-12-28}. 
\newpage
\subsection{Implementação em R}

\begin{rcode}[caption={Extração de amostra},breaklines=true, breakatwhitespace=true]
    # Instala e carrega o pacote para leitura de arquivos .csv
    install.packages("readr") 
    library(readr) 
    
    # Registra os nomes das colunas
    nomes_colunas <- c("","instant","dteday","season",
                        "weathersit","temp","casual","registered")
    
    # Extrai a amostra de 300 observações
    data_group <- read_csv("HW1_bike_sharing.csv", skip = 62, n_max = 300)
    colnames(data_group) <- nomes_colunas
    \end{rcode}


\section{Questão 2}

\section{Questão 3}

\section{Questão 4}



\end{document}